\chapter{Дрон Wing изнутри}

Дрон Wing оказался не микроконтроллером с прошивкой, а настоящим маленьким компьютером с Linux. Это была самая неожиданная находка: на борту работает та же операционная система, что на роутере или одноплатнике, а полётная логика --- обычная программа в ней.

\section{Железо}

\begin{tabularx}{\textwidth}{@{}p{3.3cm} Y@{}}
\toprule
\hdr{Узел} & \hdr{Что известно} \\ \midrule
Процессор & Allwinner семейства sun8i (класс A33), ядро Linux \texttt{3.4.39-rt143} (с патчем реального времени), имя узла \texttt{sun8i}. \pr{} \\
Память системы & Около 29 МБ оперативной памяти, из них свободно около 10 МБ при работающем автопилоте. \pr{} \\
Wi-Fi & Чип XR819 (драйверы \texttt{xradio}). В рабочем состоянии дрон выступает \emph{клиентом} сети: работают \texttt{wpa\_supplicant} и \texttt{udhcpc}. \pr{} \\
Датчики & Устройства \texttt{/dev/icm\_imu} (инерциальный модуль), \texttt{/dev/dps280} (барометр), \texttt{/dev/STM8s} (драйвер моторов на микроконтроллере STM8). \pr{} \\
Камеры & \texttt{/dev/video0} и \texttt{/dev/video1}. Вторую, по конфигурации шины, читает программа оптического потока. \vyv{} \\
UART & \texttt{ttyS0} (связь с модулем позиционирования), \texttt{ttyS2} (назначение не выяснено). \pr{} \\ \bottomrule
\end{tabularx}

\section{Флеш: восемь разделов}

С живого дрона через ADB мы сняли побитовые копии всех разделов флеша (около 8 МБ в сумме).

\begin{tabularx}{\textwidth}{@{}l r Y@{}}
\toprule
\hdr{Раздел} & \hdr{Размер} & \hdr{Назначение} \\ \midrule
\texttt{uboot} & 256 КБ & Загрузчик. \\
\texttt{boot} & 1320 КБ & Ядро Linux. \\
\texttt{rootfs} & 2416 КБ & Корневая файловая система (сжатая, только для чтения). \\
\texttt{env} & 128 КБ & Переменные загрузчика. \\
\texttt{Netprivate} & 128 КБ & Данные пользователя: параметры, Wi-Fi, файл \texttt{config.ini}. \\
\texttt{recoveryk} & 1320 КБ & Копия ядра для восстановления. \\
\texttt{recoveryf} & 2416 КБ & Копия корневой системы для восстановления. \\
\texttt{recoverye} & 128 КБ & Копия переменных загрузчика. \\ \bottomrule
\end{tabularx}

\vspace{4pt}
Корень смонтирован как \texttt{squashfs} только для чтения. Поэтому «прошивка» дрона --- это образ, который можно заменить целиком, а изменяемое состояние живёт в небольшом разделе \texttt{/netPrivate}. \pr{}

\begin{zametka}[Почему это важно]
Наличие парных разделов (\texttt{rootfs} и \texttt{recoveryf}) означает, что у дрона есть запасной вариант системы. Образ обновления от производителя нужен именно для этого механизма. Как работает переключение, мы не разбирали. \nei{}
\end{zametka}

\section{Что запущено}

С помощью списка процессов на живом дроне мы увидели, что работает по-настоящему:

\begin{tabularx}{\textwidth}{@{}p{4.4cm} Y@{}}
\toprule
\hdr{Процесс} & \hdr{Что делает} \\ \midrule
\texttt{/img\_app/pilot/pilot\_f} & Автопилот. Именно он (а не \texttt{pilot}) запущен на нашем дроне. \pr{} \\
\texttt{flow} & Считает оптический поток по картинке нижней камеры, результат отдаёт автопилоту через очередь сообщений. \pr{} \\
\texttt{wpa\_supplicant}, \texttt{udhcpc} & Подключение к Wi-Fi как клиент. \pr{} \\
\texttt{adbd} & Отладочный сервис Android Debug Bridge: даёт оболочку по USB. \pr{} \\
\texttt{pmonitor}, \texttt{sunxi\_poweroff\_detection} & Служебные: наблюдение за питанием и процессами. \pr{} \\
\texttt{getty} на консоли 115200 & Последовательная консоль. \pr{} \\ \bottomrule
\end{tabularx}

\begin{zametka}[Что значит «pilot\_f»]
В каталоге \texttt{/img\_app/pilot/} лежат два автопилота: \texttt{pilot} и \texttt{pilot\_f}. Запущен второй. Буква, по всей видимости, означает вариант для формации (режим с несколькими дронами и клиентом Wi-Fi), но прямо это нигде не написано. \vyv{}
\end{zametka}

\section{Доступ по USB}

Сервис \texttt{adbd} запускается при загрузке и показывается компьютеру как устройство \texttt{18D1:0002}. В образе есть и второй режим: диск на весь флеш (\texttt{18D1:0001}). Мы пользовались только ADB.

\begin{opasno}[Питание]
Дрон нельзя надёжно подключить только по USB: плата не поднимается. Нужен аккумулятор (при снятых пропеллерах). ADB при этом иногда отваливается через несколько минут. \pr{}
\end{opasno}

\section{Версии прошивки}

\begin{tabularx}{\textwidth}{@{}Y p{3.2cm} p{3.2cm}@{}}
\toprule
\hdr{Файл} & \hdr{На дроне} & \hdr{В образе 1.2.8} \\ \midrule
\texttt{pilot} & 206\,972 байта & 215\,184 байта \\
\texttt{pilot\_f} & 202\,880 байт & 206\,992 байта \\
Дата файлов & март 2023 & 2025 \\ \bottomrule
\end{tabularx}

\vspace{4pt}
Прошивка на дроне старше образа от производителя. Отсюда следует практический вывод: чтобы понять, что \emph{реально} работает на вашем дроне, нужно смотреть файлы с самого дрона, а не образ. \pr{}
